\newpage
\section{Conclusioni}\label{sec:conclusioni}
Il confronto sperimentale tra il metodo BFGS tollerante al rumore, BFGS classico (implementato personalmente e in SciPy) e il metodo di discesa del gradiente conferma i vantaggi pratici dei metodi quasi-Newton: i BFGS richiedono infatti molte meno iterazioni rispetto alla discesa del gradiente e, nella maggior parte dei casi, risultano più efficienti in termini di tempo e iterazioni. Tuttavia, l’introduzione di rumore nelle valutazioni della funzione e del gradiente evidenzia differenze rilevanti tra le varianti analizzate.

L’algoritmo BFGS progettato per essere tollerante al rumore si è dimostrato più robusto e affidabile rispetto a BFGS classico: grazie alle strategie di lengthening e alla line-search adattativa: mantiene l’approssimazione dell’Hessiana definita positiva, converge con maggiore probabilità e con un numero medio di iterazioni inferiore nelle prove rumorose. La versione di SciPy risulta molto efficiente in termini di tempo e di iterazioni ma mostra un tasso di convergenza più basso in presenza di rumore, suggerendo differenze nei criteri di arresto o nelle implementazioni numeriche rispetto alle implementazioni fatte personalmente. Il metodo di discesa del gradiente, pur robusto e semplice, è penalizzato dal maggior numero di iterazioni richieste e dalle difficoltà ad operare efficacemente quando il rumore è rilevante.

In sintesi, l’algoritmo BFGS tollerante al rumore mantiene i vantaggi dei quasi-Newton (rapidità e qualità della direzione di ricerca) adattandoli con successo a scenari rumorosi, rendendolo una scelta promettente quando le valutazioni sono affette da errori o approssimazioni numeriche.